\subsection{Application de sortie et conditionnement du registre}
\label{vii:sec:salida-condicionamiento}

L’inscription précédente se compose avec le lecteur de l’appareil sur
l’état généalogique complet. Le couplage conserve l’opération qui
produit la corrélation ; l’application de sortie exprime le résultat ; le
conditionnement restreint les prolongements compatibles avec ce registre.

Soient \((\Omega_S,\Sigma_S)\) et \((\Omega_M,\Sigma_M)\) les espaces
mesurables d’états généalogiques du système et de l’appareil, obtenus à partir de
leurs cylindres et registres compatibles. Un contexte de mesure s’écrit
\(\mathfrak M_a=(U_a,q_a,\mathcal C_a,m_0)\) : \(m_0\) est la
préparation de l’appareil, \(\mathcal C_a\) fixe les compatibilités et
l’horizon de prolongement, et
\[
 U_a:\Omega_S\times\Omega_M\longrightarrow\Omega_S\times\Omega_M,
 \qquad q_a:\Omega_M\longrightarrow Y_a
\]
sont, respectivement, le couplage mesurable et le lecteur mesurable vers l'espace des résultats \((Y_a,\Sigma_{Y_a})\). Dans un régime réversible, \(U_a\) est bijectif et bimesurable ; dans la réalisation de Hilbert, l'unitarité correspondante est également exigée. La sortie individuelle est la composition
\begin{equation}
 \operatorname{Out}_a(x)
 =q_a\bigl(\operatorname{pr}_M U_a(x,m_0)\bigr).
 \label{vii:eq:out-genealogico}
\end{equation}
Elle est déterminée par l’état complet, la préparation de l’appareil et
le contexte. La description statistique d’une préparation constitue
une autre lecture de l’ensemble des états compatibles.

\begin{proposition}[Mesure des résultats et restriction à la fibre]
\label{vii:prop:medida-salida}
Soit \(\Omega_a\in\Sigma_S\) l’ensemble des états compatibles avec la
préparation et le contexte, muni d’une probabilité \(\mu_a\) sur
\((\Omega_a,\Sigma_S|_{\Omega_a})\).
L’application \(\operatorname{Out}_a|_{\Omega_a}\) est mesurable et produit
la probabilité des résultats
\begin{equation}
 \nu_a(B)=\mu_a\bigl(\operatorname{Out}_a^{-1}(B)\cap\Omega_a\bigr),
 \qquad B\in\Sigma_{Y_a}.
 \label{vii:eq:medida-imagen-salida}
\end{equation}
Pour un résultat \(y\) dont le singleton est mesurable, la fibre \(\Omega_{a,y}=\{x\in\Omega_a:\operatorname{Out}_a(x)=y\}\) est mesurable. Lorsque \(\mu_a(\Omega_{a,y})>0\),
\begin{equation}
 \mu_{a,y}(E)=
 \frac{\mu_a(E\cap\Omega_{a,y})}{\mu_a(\Omega_{a,y})}
 \label{vii:eq:condicionamiento-salida}
\end{equation}
pour \(E\in\Sigma_S|_{\Omega_a}\) définit une probabilité concentrée
sur cette fibre.
\end{proposition}
\begin{proof}
L'inclusion \(x\mapsto(x,m_0)\), le couplage, la projection sur l'appareil et son lecteur sont mesurables ; leur composition l'est également. Les images réciproques conservent les réunions disjointes et l'ensemble total, de sorte que \eqref{vii:eq:medida-imagen-salida} est dénombrablement additive et de masse un. La fibre est l'image réciproque du singleton \(\{y\}\) dans \(\Omega_a\). La restriction de \(\mu_a\) à cette fibre est une mesure positive et dénombrablement additive ; la division par sa masse positive la normalise.
\end{proof}

Pour un préfixe \(s\) de profondeur \(n\), les prolongements postérieurs
au registre forment
\[
 \operatorname{Fut}_{a,y}(s)
 =\{x\in\Omega_{a,y}:x|_n=s\}.
\]
La restriction conserve le préfixe déjà parcouru et le registre du
couplage ; elle modifie la collection de prolongements admis. La sortie
\eqref{vii:eq:out-genealogico} agit sur des états complets ; la loi des
résultats \eqref{vii:eq:medida-imagen-salida} s’obtient en transportant
la mesure de préparation par cette application. La représentation par une
matrice de densité exprime le niveau
statistique de la réalisation quantique et conserve ses conditions de
correspondance avec les fibres généalogiques.
